% Demo slide showcasing the multi-language system's "mult" (multilingual)
% case: this slide's own id-slug carries the "mult" language segment
% (ZZK-mult-Multilingual-Demo), so \ifgothamSlideLanguageSet is true but
% \slideLanguage itself is "mult" -- the footer then shows whichever
% language the *compiling deck* is actually set to via \decklanguage
% (\gothamSlideLanguageFooterTag), not a fixed value. The title text
% below switches the same way, via \multilingualtext{EN}{DE}. See
% README.md's multi-language section (and LANGUAGE_NOTES.md) for the full
% design; copy from here when creating a genuinely bilingual, light
% (image-only/title-only) slide.
%% Included in decks: test-slides
%%

\setSlideVersion{1-0-0}
\begin{standardframe}{\multilingualtext{Multilingual Demo Slide}{Mehrsprachige Demo-Folie}}
	\begin{itemize}
		\item \multilingualtext{This title switches with the deck's own \texttt{\textbackslash decklanguage} setting}{Dieser Titel folgt der \texttt{\textbackslash decklanguage}-Einstellung des Decks}
		\item \multilingualtext{The footer's language tag should match the deck, not this file's own "mult" tag}{Die Sprachkennung in der Fusszeile sollte dem Deck entsprechen, nicht der "mult"-Kennung dieser Datei}
	\end{itemize}
\end{standardframe}
